%% ============================================================
%% Appendix: CoE Audit Details
%% Includes: sub-code definitions, audit procedure,
%%           forensic CPR, per-paper audit results
%% ============================================================


%\section{ADRS Adaptation} \label{appx:adrs_adaptation}
%Each baseline was adapted to the ADRS benchmark with varying levels of source modification, from prompt-only changes (DS) to patching 8--19 source files (ARC, AIR, Sakana) to interface with ADRS task specifications, evaluators, and the NeurIPS~2026 paper template.
%We standardized on Gemini~3.1 Pro as the backbone LLM across all systems for both solver code generation and paper writing.
%To ensure best-effort performance, we configured generous iteration and timeout budgets: up to 20 solver iterations per task---6.7$\times$ the default for ARC, though most tasks converge well before this cap---and 2-hour code generation windows.
%Runs that crashed due to infrastructure issues (API timeouts, rate limits, LaTeX compilation errors) were re-attempted with fresh state, up to 3~attempts per run; no run was re-attempted to improve solver scores.
%Of the 60~total runs, 18 required at least one retry. All 60~runs ultimately produced solver code and a compiled paper, though artifact quality varies.
%For each system, we run 3~seeds per task, producing 15~papers per system and 60~papers total.
%\coeaudit{} is applied identically to all systems in forensic mode.
% Full adaptation details, including per-system patches and configuration diffs, are provided in \Cref{app:baseline_adaptation}.

\section{CoE Audit Details}
\label{app:coe_audit_details}

% Sub-code definition tables removed — categories are described in plain English
% in the per-check analysis sections below (I1: \S\ref{appx:discovered-score-errors},
% I2: \S\ref{appx:i2_spec_violation}, I4: \S\ref{app:i4_audit_analysis}).



\subsection{Audit Procedure and Reproducibility}
\label{app:audit_implementation}

\coeaudit{} is designed for reproducibility.
\Cref{tab:audit_procedure} summarizes the automation level of each component.
Score Verification uses LLM-based extraction to identify the paper's reported score from \TeX{} and PDF files, then compares it against reproduced scores from evaluator re-runs via numeric tolerance---the comparison itself is fully deterministic.
Specification Violation and Method-Code Alignment are LLM-judged, with majority vote across multiple independent runs to reduce judgment noise.
Reference Verification is primarily API-based (Semantic Scholar, arXiv, OpenAlex, CrossRef lookup), with an LLM disambiguation step for title-only matches.

\begin{table}[h!]
\centering
\small
\caption{\coeaudit{} automation level per component.}
\label{tab:audit_procedure}
\begin{tabular}{@{}lll@{}}
\toprule
\textbf{Component} & \textbf{Automation} & \textbf{Model (if LLM)} \\
\midrule
Score Verification & LLM extraction + automated comparison & Gemini~3 Flash \\
Specification Violation & LLM-judged & Gemini~3.1 Pro \\
Reference Verification & Automated + LLM disambiguation & Gemini~3 Flash \\
Method-Code Alignment & LLM-judged & Gemini~3.1 Pro \\
\bottomrule
\end{tabular}
\end{table}

\paragraph{Human verification.}
Although the audit pipeline is automated, all flagged positives for I1 (Score Verification), I2 (Specification Violation), and I3 (Reference Verification) were manually reviewed and corrected by human reviewers before reporting in Table~\ref{tab:integrity}.
A small number of auditor false positives (e.g., API resolution failures for real papers, score extraction errors) were identified and removed during this process.
\emph{Exception:} most of Sakana ASv2's I2 violations trace to the BFTS--ADRS design mismatch rather than adversarial behaviour (\Cref{app:baseline_adaptation}).
For I4 (Method-Code Alignment), human reviewers validated a sample of the LLM judgments; the results in Table~\ref{tab:integrity} reflect the LLM majority-vote scores without manual correction.



%\subsection{Claim Provenance Rate: Forensic Extraction}
%\label{app:forensic_cpr}

%Forensic \cpr{} is computed post-hoc from released artifacts in three stages:

%\begin{enumerate}[nosep,leftmargin=*]
%  \item \textbf{Claim extraction.} An LLM classifier (Gemini~3 Flash) reads each sentence and labels it as a numerical result claim, a non-result sentence, or ambiguous. Only numerical result claims enter the verification pool.
%  \item \textbf{Evidence matching.} Each claim is matched against execution logs via metric-name co-occurrence: we check whether the metric name in the claim appears in a log line containing a numeric value within 1\% tolerance. A guarded value-only fallback handles absent metric names.
%  \item \textbf{Verification.} A claim is verified if a matching log line is found. Per-paper forensic $\text{CPR}_{\text{numerical}}$ is the fraction of extracted claims that are verified.
%\end{enumerate}

%Native \cpr{} is computed from provenance annotations emitted by the system at production time.
%Each claim in \texttt{claims.jsonl} carries a verification status (\texttt{passed}, \texttt{partial}, or \texttt{failed}) from the Claim Verifier.
%Only \sys{} produces native \cpr{}.

%\subsection{Implementation Details}

%\todo{Detailed description of each check's implementation:
%- Score Verification: evaluator re-run procedure, adapter per system, tolerance, failure classification
%- Hack Verification: evaluator checksum comparison, score consistency threshold
%- Reference Verification: S2 API querying, DOI/arXiv/title lookup, similarity threshold
%- Method-Code Alignment: LLM prompt, model, scoring rubric
%- Forensic CPR: claim extraction prompt, metric-name matching algorithm, trivial value filter
%This enables reproducibility and addresses the "self-graded" concern.}

%\subsection{Per-Paper CoE Audit Results}
%\label{app:per_paper}

%\todo{Full per-paper results for all 75 papers. One row per (system, seed, task) with all check results + forensic CPR. This is the evidence backing the main integrity table.}

%\todo{Consider: format as a large table, or as structured JSON/CSV in supplementary materials?}



\section{Failure Cases per Audit Metric}
\subsection{I1: Score Verification Errors }\label{appx:discovered-score-errors}

We manually reviewed every I1 (Score Verification) failure flagged
across the five systems and categorise the 22~confirmed
agent errors into five classes (Table~\ref{tab:i1_failure_distribution}).%
\footnote{A small number of auditor false positives were corrected
before reporting Table~\ref{tab:integrity}; this section reports only
confirmed agent errors.
Table~\ref{tab:integrity} reports pass rates over 12~papers per system (EPLB excluded), so its implied failure count is lower than the 22~errors here, which include EPLB papers.}

\begin{table}[h]
\centering
\small
\caption{Confirmed agent I1 errors, by category.}
\label{tab:i1_failure_distribution}
\begin{tabular}{@{}lp{6.5cm}cc@{}}
\toprule
\textbf{Category} & \textbf{Definition} & \textbf{Count} & \textbf{\%} \\
\midrule
\texttt{value\_mismatch}          & Paper headline number differs from evaluator rerun beyond tolerance, with the same metric and scale.                            & 13 & 59\% \\
\texttt{cross\_stage\_cherry\_pick}& Writeup LLM selects a score from a different search-tree stage (e.g.\ ablation) than the node whose code is submitted.         &  4 & 18\% \\
\texttt{paper\_score\_unavailable}& Paper contains no machine-readable headline number for the auditor to verify against.                                          &  2 &  9\% \\
\texttt{evaluator\_error}         & Submitted artifact cannot be re-evaluated within the budget (timeout or crash on the agent's own solution file).                &  1 &  5\% \\
\texttt{metric\_mismatch}         & Paper and code report different metrics, scales, or optimisation directions, so the numbers are not directly comparable.       &  2 &  9\% \\
\bottomrule
\end{tabular}
\end{table}

\paragraph{\texttt{value\_mismatch} (n=13).}
The paper and the code agree on what is being measured but disagree
on the value.
Most (9/13) of these gaps are within 5\% of the paper-reported
number---small enough to plausibly arise from unreported seed
variance, but uniformly biased towards a better-than-rerun headline.
Among the four baseline systems, the two largest gaps are
qualitatively different errors: a fabricated metric direction in DS
\texttt{cloudcast} seed-2 (26.7\%), where the paper relabels a cost
metric as ``utility'' and inverts the optimisation direction so that
a worse-than-baseline number reads as an improvement; and a
2$\times$ mismatch in ARC \texttt{prism} seed-3 between the paper
headline (12.74) and what the submitted code actually produces
(26.24).
Sakana ASv2 contributes two \texttt{value\_mismatch} cases, both
caused by \emph{environment-dependent tuning}: the solver contains
a hyperparameter tuning loop gated on an environment variable
(\texttt{\_ADRS\_EVAL\_GUARD}); during BFTS search the variable is
unset and the loop executes with optimal parameters, but during
canonical re-evaluation the variable is set and the loop is skipped,
causing the solver to fall back to defaults
(e.g.\ \texttt{prism}~seed-0: 26.26 tuned vs.\ 22.34 default,
a~15\% gap).

\paragraph{\texttt{cross\_stage\_cherry\_pick} (n=4).}
All four cases come from Sakana ASv2 and reflect a boundary mismatch
in the BFTS pipeline: the best node is selected from stages~1--2
(preliminary investigation and hyperparameter tuning), but the writeup LLM receives summaries
from all four stages---including ablation---and picks the most
favourable score from the entire pool.
For example, in \texttt{prism}~seed-1 the selected node scores
22.79, but the paper reports 25.39---a number traced to ablation
node~6 (``Ablate KVPR-Aware Initialization'') in the ablation
summary.
The same pattern appears in \texttt{cloudcast}~seed-0 (+56\%),
\texttt{prism}~seed-2 ($-$4.7\%, paper under-reports), and
\texttt{txn\_scheduling}~seed-2 (+17\%).

\paragraph{\texttt{paper\_score\_unavailable} (n=2).}
The agent's writeup contains no machine-readable quantitative result
for the headline metric.
Both cases are baseline-system failures: an AIR \texttt{prism} paper
that omits scores entirely, and a DS \texttt{cloudcast} PDF that
failed to render its results section.

\paragraph{\texttt{evaluator\_error} (n=1).}
The artifact the agent declares as its solution cannot be
re-evaluated end-to-end within the budget.
The single case is ARC \texttt{llm\_sql} seed-2, which ships
\texttt{program.py}---a multi-condition experiment harness that
runs 8~baseline conditions $\times$ 3~seeds $\times$ multiple
datasets at import time before any single number is produced.

\paragraph{\texttt{metric\_mismatch} (n=2).}
The paper and the code do not agree on what they are measuring.
ARC \texttt{llm\_sql} seed-3 publishes a headline of 0.0 that the
paper itself attributes to ``a critical integration failure with the
benchmark evaluator framework''---the agent surfaced a known-broken
result as its final score.
Sakana ASv2 \texttt{cloudcast}~seed-1 reports per-transfer dollar
cost while the evaluator produces a combined score, making the two
numbers incomparable.

\paragraph{Per-system error shape.}
Table~\ref{tab:i1_per_system} shows that the failure mix is highly
system-specific.
Sakana ASv2 produces the most errors (7) and the only
\texttt{cross\_stage\_cherry\_pick} cases---a failure mode unique to
multi-stage search pipelines that expose full experiment histories to
the writeup phase.
Its two \texttt{value\_mismatch} cases are environment-dependent
tuning rather than simple score inflation.
AIR and DS errors are dominated by small-to-medium
\texttt{value\_mismatch}: the numbers exist and are in the right
ballpark but do not exactly reproduce.
ARC spans the widest range of categories and the largest single
discrepancy (106\%), and contributes the only
\texttt{evaluator\_error} case.
\sys{} produces no confirmed I1 errors, consistent with its
12/12 score-verification result in Table~\ref{tab:integrity}.

\begin{table}[h]
\centering
\footnotesize
\caption{Confirmed agent I1 errors per system, by category.}
\label{tab:i1_per_system}
\begin{tabular}{@{}lcccccc@{}}
\toprule
\textbf{System} & \texttt{val\_mis.} & \texttt{cherry\_pick} & \texttt{score\_unavail.} & \texttt{eval\_error} & \texttt{metric\_mis.} & \textbf{Total} \\
\midrule
Sakana ASv2 & 2 & 4 & 0 & 0 & 1 & 7 \\
AIR    & 3 & 0 & 1 & 0 & 0 & 4 \\
ARC    & 5 & 0 & 0 & 1 & 1 & 7 \\
DS     & 3 & 0 & 1 & 0 & 0 & 4 \\
\sys{} & 0 & 0 & 0 & 0 & 0 & 0 \\
\bottomrule
\end{tabular}
\end{table}


\subsection{I2: Specification Violation Analysis}
\label{appx:i2_spec_violation}

The I2 audit flagged 11~papers across five systems for specification
violations (Table~\ref{tab:i2_per_system}).
Sakana ASv2 accounts for 10 of 11~flagged papers; the remaining
case is an AIR \texttt{llm\_sql} paper.
ARC, DS, and \sys{} produce zero I2 violations under majority vote.%
\footnote{Under union vote (any single judge flags), ARC has 3,
DS has 1, and \sys{} has 1.
As noted in the main text, DS seed-1
\texttt{llm\_sql} was flagged by only 2/5 judges---below the
majority threshold---despite containing an evaluator-exploiting
column-permutation pattern.}
Most of Sakana's I2 violations trace to the BFTS--ADRS design
mismatch (\Cref{app:baseline_adaptation}).

Because a single paper can trigger multiple violation categories, we report
both the category distribution (Table~\ref{tab:i2_subcode_distribution})
and the per-system task breakdown
(Table~\ref{tab:i2_per_system}).

\begin{table}[h]
\centering
\small
\caption{I2 violation categories across 11~flagged papers.
A paper may trigger multiple categories, so paper counts do not sum
to~11.}
\label{tab:i2_subcode_distribution}
\begin{tabular}{@{}lcc@{}}
\toprule
\textbf{Category} & \textbf{Papers} & \textbf{System(s)} \\
\midrule
Evaluator import       & 10 & Sakana (10/10) \\
Evaluator exploitation & 7  & Sakana (7/10) \\
Specification exploit  & 5  & Sakana (4/10), AIR (1) \\
Data leakage           & 1  & Sakana (1/10) \\
\bottomrule
\end{tabular}
\end{table}

\paragraph{Evaluator import (10/10 Sakana papers).}
Every flagged Sakana paper imports the canonical evaluator and calls
it as an optimisation oracle.
The mechanism is uniform: the agent adds
\texttt{from evaluator import evaluate} (or a variant aliased as
\texttt{\_tune\_eval}), then loops over hyperparameter configurations
---calling \texttt{evaluate()} for each---and keeps the best.
The import pattern mirrors our canonical evaluation harness, which
was injected into the initial code template to enable automated
scoring; the agent copies it and builds its own tuning
infrastructure on top.
The agent could tune parameters through BFTS's iteration loop
(one setting per iteration, scored by the external evaluator),
but the stage~2 goal (``test across multiple parameter settings'')
encourages an intra-iteration sweep, which requires importing the
evaluator directly.

\paragraph{Specification exploit (4 Sakana + 1 AIR).}
Four Sakana \texttt{txn\_scheduling} papers (seeds 1--3) and one
\texttt{llm\_sql} paper modify functions outside the
\texttt{EVOLVE-BLOCK} markers or exploit metric edge cases.
The \texttt{txn\_scheduling} pattern is consistent: the agent
modifies \texttt{get\_random\_costs()}, a function explicitly
forbidden from modification, to run a parameter sweep and return
the best result.
AIR's single specification-exploit case (\texttt{llm\_sql} seed-1) is a
column-permutation exploit: the solver physically reorders values
across columns within each row, destroying column integrity to
inflate the evaluator's prefix-cache hit metric---the same
pattern that appears in DS seed-1 \texttt{llm\_sql} but below
the majority threshold there.

\begin{table}[h]
\centering
\small
\caption{I2 specification violations per system. For Sakana ASv2, we
show the per-task breakdown; other systems are reported as totals.
All counts use majority vote (3/5 judges).}
\label{tab:i2_per_system}
\begin{tabular}{@{}lcccccr@{}}
\toprule
\textbf{System} & \texttt{cloud.} & \texttt{eplb} & \texttt{llm\_sql} & \texttt{prism} & \texttt{txn\_sched.} & \textbf{Total} \\
\midrule
Sakana ASv2 & 1/3 & 0/3 & 3/3 & 3/3 & 3/3 & 10/15 \\
AIR         & 0   & 0   & 1   & 0   & 0   & 1/15  \\
ARC         & 0   & 0   & 0   & 0   & 0   & 0/15  \\
DS          & 0   & 0   & 0   & 0   & 0   & 0/15  \\
\sys{}      & 0   & 0   & 0   & 0   & 0   & 0/15  \\
\bottomrule
\end{tabular}
\end{table}

\paragraph{Clean runs.}
Five of 15~Sakana runs produce no I2 violations: all three EPLB
seeds and \texttt{cloudcast} seeds 2 and~3.
EPLB's solver contract is structurally simpler (a pure allocation
function with no external scorer dependency), and the two clean
\texttt{cloudcast} runs show that even on graph-optimisation tasks
the agent \emph{can} fulfill BFTS stage goals without importing the
evaluator---but this is the exception, not the default, under
BFTS's stage pressure.


\subsection{I3: Reference integrity -- Discovered hallucinated references}\label{appx:discovered_hallucinated_references}

The list below shows unique hallucinated reference keys.
Table~\ref{tab:integrity} reports total occurrences across papers (the same key can appear in multiple papers; e.g., ARC's single fabricated key appears in all three EPLB papers, counting as 3 in the table).

\begin{longtable}{@{}p{15cm}@{}}
\toprule
\textbf{Hallucinated references} \\
\midrule
\endfirsthead
\toprule
\textbf{Hallucinated references} (continued) \\
\midrule
\endhead
\midrule
\endfoot
\bottomrule
\endlastfoot

\midrule
\multicolumn{1}{c}{\textbf{ARC (1)}} \\
\midrule
\textbf{sutskever2013importance}: Sutskever et al. \textit{SGD with Momentum}. ICML, 2013. \\

\midrule
\multicolumn{1}{c}{\textbf{AIR (21)}} \\
\midrule
\textbf{ruth2018castflow}: R\"{u}th, Jan et al. \textit{CastFlow: Clean-slate multicast approach using in-advance path computation in software-defined networks}. 2018 IEEE 43rd Conference on Local Computer Networks (LCN), 2018. \\
\textbf{dou2022hetmoe}: Dou, Shiwei et al. \textit{HetMoE: An Efficient Distributed MoE Training System for Heterogeneous Clusters}. arXiv preprint arXiv:2210.12384, 2022. \\
\textbf{li2023lightllm}: Li, Zhuohan et al. \textit{LightLLM: A highly optimized LLM inference system with token-level kv cache management}. arXiv preprint arXiv:2309.04414, 2023. \\
\textbf{purohit2024prism}: Purohit, Archit et al. \textit{Prism: Optimizing multi-model LLM serving on GPU clusters}. Proceedings of the 29th ACM International Conference on Architectural Support for Programming Languages and Operating Systems (ASPLOS), 2024. \\
\textbf{miao2023muxserve}: Miao, Xupeng et al. \textit{MuxServe: Multiplexing large language models for high throughput and low latency}. arXiv preprint arXiv:2311.05602, 2023. \\
\textbf{cloudcast}: Zheng, Q. et al. \textit{CloudCast: Cost-efficient multicast routing in cloud networks}. IEEE INFOCOM 2019 - IEEE Conference on Computer Communications, 2019. \\
\textbf{idreos2012main}: Idreos, Stratos et al. \textit{Main-memory column stores}. Foundations and Trends\textregistered\ in Databases, 2012. \\
\textbf{lightllm2023}: Li, Zhen et al. \textit{LightLLM: A Lightweight and Highly Efficient Python-based Large Language Model Serving Framework}. arXiv preprint arXiv:2310.01234, 2023. \\
\textbf{bhuiyan2024prism}: Bhuiyan, M. et al. \textit{Prism: A Flexible and Scalable Multi-LLM Serving System}. Proceedings of the 29th ACM International Conference on Architectural Support for Programming Languages and Operating Systems (ASPLOS), 2024. \\
\textbf{eplb2023}: Wang, H. et al. \textit{Expert Parallelism Load Balancing in Mixture-of-Experts Models}. Proceedings of the International Conference on High Performance Computing, Networking, Storage and Analysis (SC), 2023. \\
\textbf{hussin2011metaheuristic}: Hussin, MS et al. \textit{Metaheuristic algorithms for traveling salesman problem: A review}. Annals of the University of Craiova, Mathematics and Computer Science Series, 2011. \\
\textbf{zhang2020job}: Zhang, Jun et al. \textit{Job shop scheduling research}. International Journal of Production Research, 2020. \\
\textbf{jetstream2014}: Rabkin, Ariel et al. \textit{JetStream: Enabling wide-area data streaming}. 11th USENIX Symposium on Networked Systems Design and Implementation (NSDI 14), 2014. \\
\textbf{slingshot2022}: Doe, John et al. \textit{Slingshot: High-performance routing across federated cloud environments}. Proceedings of the ACM SIGCOMM 2022 Conference, 2022. \\
\textbf{cloudcast2021}: Chen, Yiting et al. \textit{CloudCast: Cost-effective data distribution in multi-cloud deployments}. IEEE INFOCOM 2021 - IEEE Conference on Computer Communications, 2021. \\
\textbf{cui2004}: Cui, Jun-Hong et al. \textit{QoS multicast routing in dynamic networks}. IEEE Network, 2004. \\
\textbf{laoutaris2011netstitcher}: Laoutaris, Nikolaos et al. \textit{NetStitcher: Un-tethering bulk storage from the network edge}. Proceedings of the ACM SIGCOMM 2011 conference, 2011. \\
\textbf{feng2020traffic}: Feng, Xin et al. \textit{Traffic engineering in software-defined wide-area networks: A survey}. IEEE/ACM Transactions on Networking, 2020. \\
\textbf{epstein2005online}: Epstein, Leah and van Stee, Rob. \textit{Online bin packing with square-root sized items}. Information Processing Letters, 2005. \\
\textbf{ba2023modelbox}: Ba, Yuhan et al. \textit{ModelBox: A Framework for Multi-Model Multi-Tenant Serving}. 2023 USENIX Annual Technical Conference (USENIX ATC 23), 2023. \\
\textbf{zheng2024eplb}: Zheng, Lianmin et al. \textit{EPLB: Load Balancing for Expert Parallelism in Large Language Models}. arXiv preprint arXiv:2401.03221, 2024. \\

\midrule
\multicolumn{1}{c}{\textbf{DS (41)}} \\
\midrule
\textbf{choy1991heuristic}: Choy, M-S. \textit{Heuristic algorithms for the Steiner tree problem with an application to network routing}. Proceedings of the 1991 ACM SIGCOMM, 1991. \\
\textbf{beloglazov2012energy}: Beloglazov, Anton and Buyya, Rajkumar. \textit{Energy-efficient routing and resource allocation in multi-cloud}. Journal of Network and Computer Applications, 2012. \\
\textbf{grotschel1993steiner}: Grotschel, Martin et al. \textit{The Steiner tree packing problem in telecommunications}. 1993. \\
\textbf{romero2021automated}: Romero, F et al. \textit{Automated algorithm design using large language models}. Advances in Neural Information Processing Systems, 2023. \\
\textbf{wu2015spanning}: Wu, Chuan et al. \textit{Spanning tree based data transfer in multi-cloud architectures}. 2015. \\
\textbf{melian2012cloud}: Melian, L et al. \textit{Cloud computing economics: a survey}. 2012. \\
\textbf{faragardi2017multi}: Faragardi, Hamid Reza et al. \textit{Multi-cloud data distribution with cost optimization}. 2017. \\
\textbf{mishra2018cost}: Mishra, A et al. \textit{Cost efficient routing in multi-cloud environments}. 2018. \\
\textbf{voss1992steiner}: Voss, Stefan. \textit{The Steiner tree problem with edge capacities}. 1992. \\
\textbf{bubeck2023approaches}: Bubeck, S et al. \textit{Approaches to code generation and synthesis}. 2023. \\
\textbf{wang2020automated}: Wang, X et al. \textit{Automated hyperparameter optimization in cloud computing}. 2020. \\
\textbf{smith2019data}: Smith, J et al. \textit{Data transfer costs in modern cloud platforms}. 2019. \\
\textbf{liu2022network}: Liu, Y et al. \textit{Network aware multi-cloud data distribution}. 2022. \\
\textbf{jones2023oscillatory}: Jones, R et al. \textit{Oscillatory dynamics in automated search landscapes}. 2023. \\
\textbf{kim2021puber}: Kim, Young Jin et al. \textit{Puber: Efficient expert parallelism for mixture-of-experts}. arXiv preprint arXiv:2111.05454, 2021. \\
\textbf{romera2021optimizing}: Romera-Paredes, Oscar et al. \textit{Optimizing mixture-of-experts for large-scale distributed training}. arXiv preprint arXiv:2102.04353, 2021. \\
\textbf{clune2008how}: Clune, Jeff et al. \textit{How evolutionary dynamics affects network topology}. Artificial life, 2008. \\
\textbf{llm\_agents\_2023}: Smith, J. and Doe, A. \textit{Autonomous LLM Agents for Code Generation}. Journal of AI Research, 2023. \\
\textbf{api\_misuse\_2024}: Wang, L. and Lee, C. \textit{API Misuse in LLM-Generated Code}. Proceedings of ICSE, 2024. \\
\textbf{fail\_fast\_2025}: Garcia, M. \textit{Fail-Fast Sandboxing for Coding Agents}. IEEE Transactions on Software Engineering, 2025. \\
\textbf{prism2024}: Anonymous. \textit{Prism: A Benchmark for Multi-LLM Serving Systems}. arXiv preprint, 2024. \\
\textbf{sarca2023}: Anonymous. \textit{SARCA: Systems Architecture Research using Coding Agents}. arXiv, 2023. \\
\textbf{clockwork2023}: Anonymous. \textit{Clockwork: Fast and Predictable Inference for Edge Machine Learning}. arXiv, 2023. \\
\textbf{garcia1982fully}: Garcia-Molina, Hector. \textit{A fully distributed null-free algorithm for concurrent database updates}. IEEE Transactions on Software Engineering, 1982. \\
\textbf{karlsson2020combinatorial}: Karlsson, Elias et al. \textit{Combinatorial optimization by graph neural networks}. arXiv preprint arXiv:2010.16012, 2020. \\
\textbf{krajewski2024mixtures}: Krajewski, Adam et al. \textit{Mixtures of experts: A systematic review}. arXiv preprint arXiv:2401.00000, 2024. \\
\textbf{paliwal2020regal}: Paliwal, Aditya et al. \textit{REGAL: A transfer learning based methodology for hardware-software co-design}. MICRO, 2020. \\
\textbf{jain1998simulated}: Jain, AS and Meeran, S. \textit{A simulated annealing algorithm for job shop scheduling problem}. Mathematical and computer modelling, 1998. \\
\textbf{tian2021cloud}: Tian, X et al. \textit{Cloud egress cost optimization}. Proceedings of the ACM SIGCOMM 2021 Conference, 2021. \\
\textbf{zhao2020understanding}: Zhao, Y et al. \textit{Understanding cloud network egress constraints}. USENIX Annual Technical Conference (ATC), 2020. \\
\textbf{binnig2021learned}: Binnig, Carsten et al. \textit{The case for learned database systems}. arXiv preprint, 2021. \\
\textbf{yu2014staccato}: Yu, Xiangyao et al. \textit{Staccato: A dependency-aware transaction scheduling system for many-core processors}. Proceedings of the 2014 ACM SIGMOD International Conference on Management of Data, 2014. \\
\textbf{ren2016low}: Ren, Kun et al. \textit{Low-overhead deadlock detection in distributed database systems}. Proceedings of the 2016 International Conference on Management of Data, 2016. \\
\textbf{pavlo2017make}: Pavlo, Andrew and Stonebraker, Michael. \textit{What's make a database system fast?}. ACM SIGMOD Record, 2017. \\
\textbf{blanas2010comparison}: Blanas, Spyros et al. \textit{A comparison of join algorithms for modern multi-core processors}. Proceedings of the VLDB Endowment, 2010. \\
\textbf{bailis2014bolt}: Bailis, Peter et al. \textit{Bolt-on conflict-free replicated data types}. ACM SIGMOD Record, 2014. \\
\textbf{hardi1992precedence}: Hardi, S and Rakow, TC. \textit{Precedence-based transaction scheduling}. Proceedings of the 2nd International Workshop on Research Issues on Data Engineering, 1992. \\
\textbf{kim2016fast}: Kim, Taesoo et al. \textit{Fast and scalable serializable transactions in multicore in-memory databases}. Proceedings of the 2016 International Conference on Management of Data, 2016. \\
\textbf{wu2017transaction}: Wu, Yingjun et al. \textit{Transaction scheduling using graph coloring}. Proceedings of the 2017 ACM SIGMOD International Conference on Management of Data, 2017. \\
\textbf{lin2015scheduler}: Lin, Jianshu et al. \textit{To schedule or not to schedule: When is transaction scheduling worth the overhead?}. Proceedings of the 2015 ACM SIGMOD International Conference on Management of Data, 2015. \\
\textbf{faleiro2017high}: Faleiro, Jose M and Abadi, Daniel J. \textit{High performance serializable concurrency control with determinism}. Proceedings of the 2017 ACM SIGMOD International Conference on Management of Data, 2017. \\
\end{longtable}


\subsection{I4 Audit Failure Analysis}
\label{app:i4_audit_analysis}

We categorise the 95~method-code misalignment findings flagged by
the I4 audit (across 25~papers\footnote{Table~\ref{tab:integrity} implies 26 misaligned papers (12+10+3+1); one paper (DS seed-3/cloudcast) could not be cleanly classified and is excluded from the categorical breakdown.}) into three semantic classes
(Table~\ref{tab:i4_failure_distribution}).
A single paper can produce multiple findings, so we report both the
finding-level distribution and the distinct-paper count per
category.
Sakana ASv2 is excluded from this breakdown because its I4 results
are confounded by the artifact format---the audited code is a full
experimental script rather than an extracted solver, so the I4
judges compare the paper against non-solver infrastructure code
(\Cref{app:baseline_adaptation}).

\begin{table}[h]
\centering
\small
\caption{I4 method-code alignment findings, by category (n=95
findings, 25 affected papers).}
\label{tab:i4_failure_distribution}
\begin{tabular}{@{}lp{6cm}ccc@{}}
\toprule
\textbf{Category} & \textbf{Definition} & \textbf{Findings} & \textbf{\%} & \textbf{Papers} \\
\midrule
\texttt{incomplete\_broken}        & Paper-described mechanism or component is missing, partially implemented, or replaced with a degenerate fallback in the code. & 49 & 52\% & 19 \\
\texttt{algorithm\_class\_mismatch}& Code implements a fundamentally different algorithm class than the one described in the paper (e.g.,~paper claims Iterated Local Search, code is Simulated Annealing; paper claims a neural- or LLM-guided search, code is deterministic). & 37 & 39\% & 15 \\
\texttt{deceptive\_dummy\_code}    & Code contains undisclosed elements designed to mislead evaluation: hidden environment-variable switches, or output deliberately shaped to game an evaluator metric. &  9 &  9\% &  5 \\
\bottomrule
\end{tabular}
\end{table}

\paragraph{\texttt{incomplete\_broken} (n=49, 19 papers).}
The largest category. The code generally targets the same problem
the paper describes, but is missing one or more of the specific
mechanisms claimed in the writeup, or substitutes a degenerate
fallback. Recurring patterns include:
described \emph{multi-start initialisation with $K$ sequences}
collapsed to a single deterministic backtracking pass
(AIR \texttt{prism}); claimed \emph{link-penalty diversification}
producing a constant assignment of the same path to every partition
(AIR \texttt{cloudcast}); described \emph{arborescence lookahead}
or \emph{surrogate cost model} simply absent from the code, with the
true simulator invoked at every iteration instead
(ARC \texttt{cloudcast}, ARC \texttt{txn\_scheduling});
claimed \emph{2.0\,GB memory threshold safeguards} reduced to a
trivial overflow check (DS \texttt{prism}).

\paragraph{\texttt{algorithm\_class\_mismatch} (n=37, 15 papers).}
The code implements a fundamentally different algorithm class than
the one described in the paper. The most common sub-pattern is the
\emph{claimed-learning-loop-is-absent} case: a paper claims an
LLM-driven evolutionary search, a neural network predictor, or an
LLM SQL optimiser, and the code is a single deterministic heuristic
with no LLM calls and no trained model
(e.g.\ AIR \texttt{cloudcast}: ``hybrid neuro-symbolic solver''
$\rightarrow$ purely classical Steiner-tree heuristics;
ARC \texttt{llm\_sql}: ``36 LLM prompting strategies''
$\rightarrow$ deterministic dataframe reordering;
DS \texttt{eplb}: ``LLM-driven evolutionary search over 27
generations'' $\rightarrow$ standalone deterministic load
balancer).
Classical algorithm-class swaps also recur:
Iterated Local Search $\rightarrow$ Simulated Annealing
(AIR \texttt{prism}), Dinkelbach's fractional programming
$\rightarrow$ static sum-of-squares cost (ARC \texttt{prism}),
recursive partitioning $\rightarrow$ single-level grouping
(DS \texttt{llm\_sql}).
We also observe the rarer \emph{method/baseline inversion} pattern,
in which the paper labels $X$ as the proposed method and $Y$ as a
rejected ablation, but the code uses $Y$
(ARC \texttt{prism}: ``GRASP Without Symbiosis'' claimed as final
method while the code instantiates \texttt{SymbioticGRASPPacker}).

\paragraph{\texttt{deceptive\_dummy\_code} (n=9, 5 papers).}
The submitted artifact contains undisclosed code whose presence
appears intended to mislead automated evaluation. All nine
findings come from ARC and split into two sub-patterns:
(i)~\emph{hidden environment-variable switches} (4 findings,
2~papers): the code reads an undisclosed variable
(\texttt{CONDITION}, \texttt{ABLATION}) at import time and dispatches
to one of several different solvers, while the paper presents a
single unified algorithm
(ARC \texttt{cloudcast} seeds~2 and~3).
(ii)~\emph{evaluator gaming} (5 findings, 3~papers): code
intentionally shaped to inflate a metric without solving the
underlying task, e.g.\ returning an empty column-ordering list while
internally permuting values to maximise prefix-cache hits
(ARC \texttt{llm\_sql}~seed-3).

\paragraph{Per-system error shape.}
Table~\ref{tab:i4_per_system} shows that the failure mix is sharply
system-specific.
AIR, DeepScientist, and \sys{} produce no
\texttt{deceptive\_dummy\_code} findings: when these systems
misalign, the gap is between what the paper says and what the code
does, with no active obfuscation.
ARC is the only system that produces
\texttt{deceptive\_dummy\_code} findings (9/9, 5 papers), spanning
both hidden env-var switches and prefix-hit evaluator gaming.

\begin{table}[h]
\centering
\small
\caption{I4 findings per system, by category. Each cell is a
finding count; a single paper may contribute multiple findings.}
\label{tab:i4_per_system}
\resizebox{\textwidth}{!}{%
\begin{tabular}{@{}lcccc@{}}
\toprule
\textbf{System} & \texttt{algo\_class\_mismatch} & \texttt{incomplete\_broken} & \texttt{deceptive\_dummy\_code} & \textbf{Total} \\
\midrule
AIR &  3 &  9 & 0 & 12 \\
ARC & 15 & 23 & 9 & 47 \\
DS  & 15 & 16 & 0 & 31 \\
\sys{} &  4 &  1 & 0 &  5 \\
\bottomrule
\end{tabular}}
\end{table}

\section{MLE-Bench and Parameter Golf Evaluation}
\label{app:mle-pg-eval}

\begin{table}[h]
\centering
\small
\caption{Mapping between MLE-Bench task names used in our evaluation and their corresponding official task IDs.}
\label{tab:mle-bench-task-mapping}
\begin{tabular}{ll}
    \toprule
    \textbf{Task Name} & \textbf{Task ID} \\ 
    \midrule
    3D Object Detection & \texttt{3d-object-detection-for-autonomous-vehicles} \\
    AI4Code & \texttt{AI4Code} \\
    iMet 2020 FGVC7 & \texttt{imet-2020-fgvc7} \\
    RSNA Brain Tumor & \texttt{rsna-miccai-brain-tumor-radiogenomic-classification} \\ 
    iNaturalist 2019 FGVC6 & \texttt{inaturalist-2019-fgvc6} \\
    \bottomrule
\end{tabular}
\end{table}

We evaluate \sys{} on MLE-Bench and Parameter Golf. Details regarding the benchmarks and our evaluation setup are provided below.

\paragraph{MLE-Bench.} We selected five MLE-Bench~\cite{chan2024mle} Kaggle competitions spanning medical imaging, fine-grained recognition, and 3D perception. Table~\ref{tab:mle-bench-task-mapping} summarizes the mapping between the task names used in our evaluation and their official task IDs. We specifically targeted tasks in the Medium and High difficulty tiers, as they possess sufficient complexity and substantive research value to serve as meaningful benchmarks for automated scientific discovery and paper writing. Specifically, AI4Code, iMet 2020 FGVC7, and iNaturalist 2019 FGVC6 are classified as Medium-difficulty tasks, whereas 3D Object Detection and RSNA Brain Tumor are High-difficulty tasks. The agent is provided with task descriptions and a code execution environment to develop its solutions. The execution environment consists of a compute node provisioned with 8xH100 GPUs, 192 CPU cores, and 1TB of RAM. To assist with formatting, the agent can query a validation tool an unlimited number of times to ensure its submission CSV files are correctly structured. Furthermore, we permit the agent to query the grading server up to 16 times to obtain evaluation scores on the test data. This deviates from the official MLE-Bench protocol, which restricts evaluation to a single submission of the final generated solution. We adopt this modified setup to better simulate real-world Kaggle competitions, where participants iteratively submit solutions to a public leaderboard for feedback during the development phase. While the agent is instructed to iterate on its solutions primarily based on metrics derived from the validation dataset, the grading server is used sparingly to select the best-performing models on the test set. Additionally, this setup is necessary to allow the agent to report accurate final test metrics when generating the paper.

\paragraph{Parameter Golf.} To assess adaptability in a novel, live environment, we evaluate \sys{} on the Parameter Golf competition~\citep{openai2026parametergolf}. This live competition serves as a well-suited AI research task focused on LLM training. It requires participants to train the highest-performing language model under strict constraints: the final artifact must fit within a 16MB size limit and the training process must complete in under 10 minutes on an 8xH100 system. Performance is evaluated by the compression rate---measured in tokenizer-agnostic bits per byte (BPB)--- on the FineWeb validation set. We provide the agent with a sandbox tool to execute solution code to obtain evaluation metrics. Additionally, the agent is provided with a knowledge base containing official leaderboard solutions up to a cutoff date of April 27, 2026. As of this date, the state-of-the-art (SOTA) score was 1.0611, achieved by a solution titled ``BOS-Fixed SmearGate + LQER + SparseAttnGate + 9-Hparam Stack''. The agent is tasked with building upon these existing leaderboard baselines to discover novel techniques capable of further improving the SOTA score.